\documentclass[11pt]{article}
\usepackage[margin=1in]{geometry}
\usepackage[T1]{fontenc}
\usepackage[utf8]{inputenc}
\usepackage{lmodern}
\usepackage{amsmath,amssymb}
\usepackage{hyperref}
\usepackage{xcolor}
\usepackage{listings}

\lstdefinestyle{pycode}{
  language=Python,
  basicstyle=\ttfamily\small,
  keywordstyle=\color{blue!70!black},
  commentstyle=\color{green!40!black},
  stringstyle=\color{orange!60!black},
  showstringspaces=false,
  breaklines=true,
  frame=single,
  tabsize=4
}

\title{CTM Neural Synchronization Change: Exponential Decay $\rightarrow$ RPC Filter}
\author{Implementation Note}
\date{\today}

\begin{document}
\maketitle

\section{Scope}
This document records the exact proposed and implemented change to neural synchronization in CTM:
\begin{itemize}
  \item Original mechanism: single-timescale learnable exponential decay recursion.
  \item New option: \textbf{RPC} (Robust Persistent Coherence) synchronization filter.
\end{itemize}
The change is implemented with backward compatibility:
\texttt{sync\_mode='exp'} remains the default, and \texttt{sync\_mode='rpc'} enables the new filter.

\section{Rationale}
The original update tracks weighted pairwise products, but a single exponential kernel has limitations for robustness:
\begin{itemize}
  \item Strong short-term spikes can dominate accumulated synchronization.
  \item Short coincidences are not explicitly separated from persistent structure.
  \item Weak but consistent patterns can be overshadowed by high-amplitude noise.
\end{itemize}

RPC addresses this with four mechanisms:
\begin{enumerate}
  \item \textbf{Robust product clipping}: bounds one-step outlier impact.
  \item \textbf{Dual timescales} (slow and fast): distinguishes persistent signal from transient coincidence.
  \item \textbf{Persistence gate}: retains contributions only when slow coherence exceeds fast coincidence.
  \item \textbf{Power normalization}: outputs scale-normalized coherence, reducing amplitude confounding.
\end{enumerate}

\section{Mathematical Form}
For each selected neuron pair $(i,j)$ at tick $t$:
\begin{align}
\tilde{u}_t &= \delta\,\tanh\left(\frac{(z_{i,t}-\mu_{i,t})(z_{j,t}-\mu_{j,t})}{\delta}\right), \\
 c_t^s &= \rho_s c_{t-1}^s + (1-\rho_s)\tilde{u}_t, \\
 c_t^f &= \rho_f c_{t-1}^f + (1-\rho_f)\tilde{u}_t, \quad \rho_f < \rho_s, \\
 g_t &= \sigma\big(k(|c_t^s|-|c_t^f|-\tau)\big), \\
 \mathrm{sync}_t &= g_t\,\frac{c_t^s}{\sqrt{p_{i,t}p_{j,t}+\epsilon}}.
\end{align}
Implementation maps these to learnable parameters per synchronization stream (\texttt{out}, \texttt{action}):
\texttt{rpc\_clip\_delta}, \texttt{rpc\_gate\_scale}, \texttt{rpc\_tau}, \texttt{rpc\_fast\_offset}.

\section{Exact Code Snippets (Before $\rightarrow$ After)}

\subsection{Core synchronization recursion}
\textbf{Before} (single exponential recursion in \texttt{models/ctm.py}, pre-change):
\begin{lstlisting}[style=pycode]
# Compute synchronisation recurrently
if decay_alpha is None or decay_beta is None:
    decay_alpha = pairwise_product
    decay_beta = torch.ones_like(pairwise_product)
else:
    decay_alpha = r * decay_alpha + pairwise_product
    decay_beta = r * decay_beta + 1

synchronisation = decay_alpha / (torch.sqrt(decay_beta))
return synchronisation, decay_alpha, decay_beta
\end{lstlisting}

\textbf{After} (mode-switch in \texttt{/Users/aryangoyal/Sakana AI/continuous-thought-machines\_new/models/ctm.py}):
\begin{lstlisting}[style=pycode]
pairwise_product, pairwise_left, pairwise_right = self._get_pairwise_activity(activated_state, synch_type)

if self.sync_mode == 'rpc':
    synchronisation, state = self._compute_rpc_synchronisation(
        pairwise_left=pairwise_left,
        pairwise_right=pairwise_right,
        rho_s=r,
        state=decay_alpha if isinstance(decay_alpha, dict) else None,
        synch_type=synch_type,
    )
    return synchronisation, state, None

if decay_alpha is None or decay_beta is None:
    decay_alpha = pairwise_product
    decay_beta = torch.ones_like(pairwise_product)
else:
    decay_alpha = r * decay_alpha + pairwise_product
    decay_beta = r * decay_beta + 1

synchronisation = decay_alpha / (torch.sqrt(decay_beta))
return synchronisation, decay_alpha, decay_beta
\end{lstlisting}

\subsection{RPC implementation block}
\textbf{After} (new function in \texttt{/Users/aryangoyal/Sakana AI/continuous-thought-machines\_new/models/ctm.py}):
\begin{lstlisting}[style=pycode]
def _compute_rpc_synchronisation(self, pairwise_left, pairwise_right, rho_s, state, synch_type):
    delta = F.softplus(getattr(self, f'rpc_clip_delta_{synch_type}')) + 1e-6
    gate_scale = F.softplus(getattr(self, f'rpc_gate_scale_{synch_type}'))
    tau = getattr(self, f'rpc_tau_{synch_type}')
    fast_offset = F.softplus(getattr(self, f'rpc_fast_offset_{synch_type}'))

    rho_s = torch.clamp(rho_s, min=1e-6, max=1 - 1e-6)
    rho_f = torch.clamp(rho_s * torch.exp(-fast_offset), min=1e-6, max=1 - 1e-6)
    one_minus_rho_s = 1 - rho_s
    one_minus_rho_f = 1 - rho_f

    if state is None:
        m_left = pairwise_left
        m_right = pairwise_right
        c_s = pairwise_left * pairwise_right
        c_f = c_s.clone()
        var_left = torch.ones_like(pairwise_left)
        var_right = torch.ones_like(pairwise_right)
    else:
        m_left = state['m_left']
        m_right = state['m_right']
        c_s = state['c_s']
        c_f = state['c_f']
        var_left = state['var_left']
        var_right = state['var_right']

    m_left = rho_s * m_left + one_minus_rho_s * pairwise_left
    m_right = rho_s * m_right + one_minus_rho_s * pairwise_right

    centered_left = pairwise_left - m_left
    centered_right = pairwise_right - m_right

    robust_prod = delta * torch.tanh((centered_left * centered_right) / delta)
    c_s = rho_s * c_s + one_minus_rho_s * robust_prod
    c_f = rho_f * c_f + one_minus_rho_f * robust_prod
    var_left = rho_s * var_left + one_minus_rho_s * centered_left.square()
    var_right = rho_s * var_right + one_minus_rho_s * centered_right.square()

    gate = torch.sigmoid(gate_scale * (c_s.abs() - c_f.abs() - tau))
    synchronisation = gate * c_s / torch.sqrt(var_left * var_right + 1e-6)
    new_state = {
        'm_left': m_left,
        'm_right': m_right,
        'c_s': c_s,
        'c_f': c_f,
        'var_left': var_left,
        'var_right': var_right,
    }
    return synchronisation, new_state
\end{lstlisting}

\subsection{Parameter registration (new RPC learnables)}
\textbf{Before}:
\begin{lstlisting}[style=pycode]
self.register_parameter(f'decay_params_{synch_type}', nn.Parameter(torch.zeros(synch_representation_size), requires_grad=True))
\end{lstlisting}

\textbf{After}:
\begin{lstlisting}[style=pycode]
self.register_parameter(f'decay_params_{synch_type}', nn.Parameter(torch.zeros(synch_representation_size), requires_grad=True))
if self.sync_mode == 'rpc':
    self.register_parameter(
        f'rpc_clip_delta_{synch_type}',
        nn.Parameter(torch.tensor(float(self.rpc_clip_delta_init), dtype=torch.float32), requires_grad=True),
    )
    self.register_parameter(
        f'rpc_gate_scale_{synch_type}',
        nn.Parameter(torch.tensor(float(self.rpc_gate_scale_init), dtype=torch.float32), requires_grad=True),
    )
    self.register_parameter(
        f'rpc_tau_{synch_type}',
        nn.Parameter(torch.tensor(float(self.rpc_tau_init), dtype=torch.float32), requires_grad=True),
    )
    self.register_parameter(
        f'rpc_fast_offset_{synch_type}',
        nn.Parameter(torch.tensor(float(self.rpc_fast_offset_init), dtype=torch.float32), requires_grad=True),
    )
\end{lstlisting}

\subsection{Decay setup cleanup (removed in-forward \texttt{.data} clamp)}
\textbf{Before}:
\begin{lstlisting}[style=pycode]
self.decay_params_action.data = torch.clamp(self.decay_params_action, 0, 15)
self.decay_params_out.data = torch.clamp(self.decay_params_out, 0, 15)
r_action, r_out = torch.exp(-self.decay_params_action).unsqueeze(0).repeat(B, 1), torch.exp(-self.decay_params_out).unsqueeze(0).repeat(B, 1)
\end{lstlisting}

\textbf{After}:
\begin{lstlisting}[style=pycode]
r_action = self._get_decay_factors(self.decay_params_action, B)
r_out = self._get_decay_factors(self.decay_params_out, B)
\end{lstlisting}

\section{Training Surface Changes (for reproducible finetuning)}
The training scripts were extended to support controlled experiments with this change:
\begin{itemize}
  \item \texttt{--sync\_mode exp|rpc} and RPC hyperparameters.
  \item \texttt{--init\_checkpoint\_path}/\texttt{--init\_checkpoint\_dir} for external checkpoint initialization.
  \item \texttt{--finetune\_sync\_only} to isolate synchronization-related learning.
\end{itemize}
Implemented in:
\begin{itemize}
  \item \texttt{/Users/aryangoyal/Sakana AI/continuous-thought-machines\_new/tasks/mazes/train.py}
  \item \texttt{/Users/aryangoyal/Sakana AI/continuous-thought-machines\_new/tasks/image\_classification/train.py}
\end{itemize}
For ImageNet subset reproducibility, local subset support and downloader were added:
\begin{itemize}
  \item \texttt{/Users/aryangoyal/Sakana AI/continuous-thought-machines\_new/data/custom\_datasets.py}
  \item \texttt{/Users/aryangoyal/Sakana AI/continuous-thought-machines\_new/tasks/image\_classification/download\_imagenet\_subset.py}
\end{itemize}

\section{Expected Behavior}
Compared to the exponential baseline, RPC is expected to:
\begin{itemize}
  \item Reduce sensitivity to one-step activation bursts (noise suppression).
  \item Suppress short-lived coincidences via fast/slow contrast gate.
  \item Retain weak but persistent synchrony patterns over longer horizons.
  \item Preserve backward compatibility by keeping \texttt{exp} as the default mode.
\end{itemize}

\end{document}
